We now turn to the other requirement of Freedman's inequality:
a bound on the absolute size of the increments $D_t(\bar{x})$.
An almost sure bound on the $|D_t(\bar{x})|$ is not possible, since both $Y_t$ and $G_t(\bar{x})$ are unbounded.
However, the Gaussian score is already controlled on the event $\mathcal E_G^{(s:u)}(\bar{x})$.
We do the same trick of introducing a high-probability event to control the feedback $Y_t$.
Recall that
$Y_t=\pi_\eps(X_t)(f(A_t)+\xi_t)+2v(X_t)/\eps$. Since
$f(A_t)\in[0,1]$, we have
\begin{align}
    |\eps Y_t|
    \leq\eps\pi_\eps(X_t)(1+|\xi_t|)+2v(X_t).
    \label{eq:feedback-cutoff-decomposition}
\end{align}
Thus we need to control both the observation noise and how far the sample
$X_t$ lies outside $\K_\eps$.
For the noise, define
\begin{align}
    \mathcal E_{\mathrm{noise}}
    =\left\{|\xi_t|\leq\sqrt\ell
    \text{ for all }1\leq t\leq n\right\}.
    \label{eq:noise-cutoff-event}
\end{align}
By \eqref{eq:noise}, $\Pp_t(|\xi_t|>\sqrt\ell)\leq2e^{-\ell}$,
so a union bound gives $\Pp(\overline{\mathcal E}_{\mathrm{noise}})\leq2ne^{-\ell}$.
For the sample, write $W_t=X_t-\mu_t$. Subadditivity of the gauge gives
\begin{align}
    v(X_t)\leq v(\mu_t)+\frac{\pi(W_t)}{1-\eps}
    \leq v(\mu_t)+2\pi(W_t).
\end{align}
The shell condition already bounds $v(\mu_t)$, so it remains to control
the Gaussian displacement $W_t$.
To this end, define
\begin{align}
    \mathcal E_{\mathrm{gauge}}
    =\left\{\pi(W_t)\leq3+2\sqrt\ell
    \text{ for all }1\leq t\leq\tau_{\mathrm{cov}}\right\}\,,
    \label{eq:gauge-cutoff-event}
\end{align}
which we will show holds with high probability as $\pi$ is 1-Lipschitz and $W_t$ is conditionally Gaussian
with covariance $\widehat\Sigma_t\preceq2\sigma^2I_d$.
On $\mathcal E_{\mathrm{gauge}}$, the shell condition now gives
\begin{align}
    v(X_t)\leq\ell+6+4\sqrt\ell\leq2\ell,
    \qquad\pi_\eps(X_t)=1+v(X_t)\leq3\ell
    \qquad(t\leq\tau_{\mathrm{cov}}\wedge\tau_{\mathrm{shell}}).
\end{align}
If $\mathcal E_{\mathrm{noise}}$ also holds, substituting these bounds
into \eqref{eq:feedback-cutoff-decomposition} yields
\begin{align}
    |\eps Y_t|
    \leq\frac32\ell(1+\sqrt\ell)+4\ell
    =\frac{11}{2}\ell+\frac32\ell^{3/2}
    \leq4\ell^{3/2},
    \label{eq:feedback-cutoff-bound}
\end{align}
where the numerical inequalities use $\ell\geq640$.
Now, define the truncated quadratics
\begin{align}
    \widetilde q_t(\bar{x})
    =\widehat q_t(\bar{x})\1_{\mathcal C_t(\bar{x})}\,
    \quad\text{where}\quad
    \mathcal C_t(\bar{x})
    =\left\{|\xi_t|\leq\sqrt\ell\,,\,
    \pi(W_t)\leq3+2\sqrt\ell\right\}
    \cap\mathcal E_G^{(t:t)}(\bar{x})\,,
    \label{eq:quadratic-cutoff-event}
\end{align}
which, unlike the original quadratics $\widehat q_t(\bar{x})$, satisfies the absolute bound
\begin{align}
    |\widetilde q_t(\bar{x})|
    \leq2|Y_t|R_t(\bar{x})(G_t(\bar{x})^2+1)
    \1_{\mathcal C_t(\bar{x})}
    \leq\frac{40\sqrt{2F}\,\ell^{5/2}}{\eps}.
    \label{eq:truncated-quadratic-bound}
\end{align}
To obtain martingale differences, we therefore center it at their own conditional mean:
\begin{align}
    \widetilde D_t(\bar{x})
    =\E_t[\widetilde q_t(\bar{x})]-\widetilde q_t(\bar{x}).
    \label{eq:truncated-martingale-difference}
\end{align}
By construction, $\E_t[\widetilde D_t(\bar{x})]=0$. Both the truncated
estimate and its conditional mean satisfy the preceding absolute bound,
so $|\widetilde D_t(\bar{x})| \leq\frac{80\sqrt{2F}\,\ell^{5/2}}{\eps}$.
Quite similarly to \eqref{eq:martingale-variance}, the accumulated conditional variance of the truncated martingale differences is bounded by the sum of the conditional second moments of the original quadratics as
\begin{align}
    \sum_{t=s}^{u}\E_t[\widetilde D_t(\bar{x})^2]
    \leq\sum_{t=s}^{u}\E_t[\widetilde q_t(\bar{x})^2]
    \leq\sum_{t=s}^{u}\E_t[\widehat q_t(\bar{x})^2]
    \leq200F\ell^2\sum_{t=s}^{u}\E_t[Y_t^2]
    +\frac{8F(u-s+1)}{\eps^2}e^{-\ell/4},
    \label{eq:truncated-variance-comparison}
\end{align}
where the last inequality follows from \eqref{eq:quadratic-conditional-variance}.

We can now apply Freedman's inequality, leaving the feedback second-moment
budget to be specified in the global and local arguments. Fix deterministic
$F,S>0$. Let $\mathcal E_{\mathrm{Fr}}^{(s:u)}(\bar{x};S)$ be the event that,
whenever the preceding stopping and radius conditions hold throughout
$[s,u]$ and $\sum_{t=s}^{u}\E_t[Y_t^2]\leq S$, we have
\begin{align}
    \left|\sum_{t=s}^{u}\widetilde D_t(\bar{x})\right|
    \leq\sqrt{25F\ell^3S+\frac{F(u-s+1)\ell}{\eps^2}e^{-\ell/4}}
    +\frac{10\sqrt{2F}\,\ell^{7/2}}{3\eps}.
    \label{eq:freedman-quadratic-bound}
\end{align}

It remains to relate this martingale sum to the original error. On
$\mathcal E_{\mathrm{Fr}}^{(s:u)}(\bar{x};S)
    \cap\mathcal E_{\mathrm{noise}}\cap\mathcal E_{\mathrm{gauge}}
    \cap\mathcal E_G^{(s:u)}(\bar{x})$, every cutoff $\mathcal C_t(\bar{x})$
holds for $s\leq t\leq u\leq\tau_{\mathrm{cov}}$.
The sampled and truncated quadratics therefore agree on the interval.
Their conditional means still differ by the expected contribution of
the discarded samples. Summing these differences gives
\begin{align}
    \left|\sum_{t=s}^{u}D_t(\bar{x})\right|
    \leq\left|\sum_{t=s}^{u}\widetilde D_t(\bar{x})\right|
    +\sum_{t=s}^{u}\left|
    \E_t\left[\widehat q_t(\bar{x})
        \1_{\overline{\mathcal C_t(\bar{x})}}\right]\right|.
    \label{eq:truncation-bias-decomposition}
\end{align}
On this intersection, under the preceding stopping and radius conditions
and the feedback budget above, the first term is bounded by
\eqref{eq:freedman-quadratic-bound}. It remains to bound the second term,
the bias introduced by truncation.



% To control this difference, we have to rely on a couple of high-probability events.
% First, let
% \begin{align}
%     \mathcal{E}_{\mathrm{noise}} = \left\{...\right\}
%     \quad\text{and}\quad
%     \mathcal{E}_{\mathrm{gauge}} = \left\{...\right\}
% \end{align}
% where $\mathcal{E}_{\mathrm{noise}}$ is the event that the noise is well-behaved and $\mathcal{E}_{\mathrm{gauge}}$ is the event that 


